\documentclass[12pt,letterpaper]{article}
\usepackage[T1]{fontenc}
\usepackage{mathptmx}
\usepackage{courier}
\usepackage[letterpaper,margin=1in]{geometry}
\usepackage{amsmath,amssymb}
\usepackage{microtype}
\usepackage{booktabs}
\usepackage{array}
\usepackage{longtable}
\usepackage{enumitem}
\usepackage{xr}
\externaldocument[N-]{infres}
\usepackage[hidelinks]{hyperref}

\newcommand{\C}{\mathbb{C}}
\newcommand{\Z}{\mathbb{Z}}
\newcommand{\act}[2]{{}^{#1}\!#2}
\DeclareMathOperator{\Hom}{Hom}
\DeclareMathOperator{\Aut}{Aut}
\DeclareMathOperator{\Int}{Int}
\DeclareMathOperator{\SL}{SL}
\DeclareMathOperator{\Cay}{Cay}

\title{Revision log\\[4pt]\large ``Notes on the inflation-restriction
sequence in non-abelian cohomology'', Draft 2\\(corrected and expanded
version)}
\author{David Joyner\\\texttt{wdjoyner@gmail.com} \and Claude}
\date{November 26, 2026}

\begin{document}
\maketitle

\noindent This log records every change from the original manuscript to
Draft~1 and from Draft~1 to Draft~2, the responses to the two referee
reports, the AI usage, and the sources consulted and not consulted.
Numbers of statements are those of the current draft (Draft~2) unless a
draft is named; they are generated from the note's \texttt{.aux} file, so
compile the note first. Section~\ref{sec:d1d2} has a concordance with the
numbering of Draft~1.

\section{Files and drafts}\label{sec:files}

\begin{itemize}[itemsep=1pt]
\item All files are in the folder \texttt{infres} of the repository
  \url{https://github.com/wdjoyner/alg-graph-thry}.
\item Original: \texttt{original/infres.tex}, an undated, unpublished
  \AmS-\TeX\ manuscript of about four pages, written while the author was
  at the U.S.\ Naval Academy. It is kept unchanged. It needs plain \TeX\
  with \texttt{amstex} and \texttt{amsppt} to compile.
\item Draft 1 (November 26, 2026). It exists in two states. The state
  first pushed (commit \texttt{204c988}) was merged into the main branch
  as pull request~\#14; it lists D.~Joyner as sole author and has a short
  glossary and no running example. The amended state
  (Section~\ref{sec:amend}; commits \texttt{2336a5f} and \texttt{d57e53e})
  was pushed to the same branch after the merge and so did not reach the
  main branch until Draft~2.
\item Draft 2 (November 26, 2026): \texttt{infres.tex} and its PDF. It
  starts from the amended Draft~1.
\item Referee report on the original, by Gemini (``Carbon''), as
  received: the file \nolinkurl{carbon_referee_report.md} in the folder
  \texttt{referee/2026-11-26/}.
\item Referee report on Draft~1, by ChatGPT, as received: the file
  \nolinkurl{chatgpt_referee_report.md} in the folder
  \texttt{referee/2026-11-26-chatgpt/}. It was written from the state of
  Draft~1 on the main branch, that is, the state first pushed.
\item Script: \texttt{verification/infres\_checks.py} (Python~3, standard
  library only), with its output in
  \texttt{verification/checks\_output.txt}. A full run takes about six
  minutes.
\end{itemize}

\section{AI usage}\label{sec:ai}

AI review is not an independent certification of the results.

\begin{enumerate}[itemsep=1pt]
\item \emph{Referee.} Gemini (Google), called ``Carbon'' by the author,
  refereed the original manuscript. Its report is reproduced in
  \texttt{referee/2026-11-26/}.
\item \emph{Reviser and author of new material.} Claude (Anthropic, model
  Claude Fable~5.2) prepared Draft~1 at the author's request, with the
  instructions: revise according to the referee report, apply the house
  style sheet, fill in details, look for graph-theoretic connections that
  organize the point of view, and post the result to the repository.
  Claude checked each referee item independently
  (Section~\ref{sec:referee}), wrote all new text, proofs and examples, the
  literature paragraph, the glossary and Sections~\ref{N-sec:graph}
  and~\ref{N-sec:langlands} of the note.
\item \emph{Computation.} Claude wrote and ran the verification script
  (Section~\ref{sec:verification}).
\item \emph{Literature.} Claude ran one web search to check how the
  non-abelian inflation-restriction sequence and Springer's $H^2$ are cited
  in recent papers (Section~\ref{sec:sources}).
\end{enumerate}

\paragraph{Draft 2 (November 26, 2026).}
\begin{enumerate}[itemsep=1pt]
\item \emph{Referee.} ChatGPT (OpenAI) refereed Draft~1 and recommended
  minor to moderate revision. Besides corrections it proposed new
  mathematics with proofs or proof sketches
  (Section~\ref{sec:chatgpt}).
\item \emph{Reviser and second author.} Claude (Anthropic, model Claude
  Fable~5.2) checked every item of the report, wrote out and verified the
  proposed results, added further results and examples, and extended the
  verification script (Section~\ref{sec:d1d2}).
\item No web search was made for Draft~2.
\end{enumerate}

\noindent As first pushed, Draft~1 listed D.~Joyner as sole author. At
his request (November 26, 2026) Claude was added as second author; see
Section~\ref{sec:amend}.

\section{Response to the referee report}\label{sec:referee}

Items are numbered in the order of the report.

\begin{longtable}{@{}p{0.05\textwidth}p{0.40\textwidth}p{0.49\textwidth}@{}}
\toprule
No. & Referee's point & Response \\
\midrule
\endhead
1.1 & $E^\times$ is $W_E/W_E^c$, not $W_F/W_F^c$. & Accepted. See the
first paragraph of the Introduction and (W1). \\
1.2 & The hypothesis $s(w^{-1}gw)=w^{-1}s(g)w$ on the section is
ill-typed and too strong; require only $p(s(g))=g$, $s(1)=1$. The class of
$\sigma$ is non-trivial only if $E\neq F$. & Accepted.
Definition~\ref{N-def:section}; the hypothesis $E\neq F$ is stated in
\S\ref{N-ss:scope} and after (W3). \\
1.3 & $H^1(E^\times,A)^\Gamma$ is not
$\Hom_\Gamma(E^\times,A)/A$: an invariant class need not contain an
equivariant homomorphism, and conjugates of equivariant homomorphisms need
not be equivariant. & Accepted. Definition~\ref{N-def:T} separates
``invariant class'' from ``equivariant homomorphism'', and
Example~\ref{N-ex:toy-noeq} and Example~\ref{N-ex:dihedral} give invariant
classes with no equivariant representative. \\
1.4 & Fix: keep Springer's full definition, with elements $b_g$ and a
twisted action; define the $W$-action on $\Hom(E^\times,A)$ by
$(w\star\phi)(x)=\act{w}{\phi(w^{-1}xw)}$; the restriction of a cocycle
satisfies $f(wxw^{-1})=f(w)\,\act{w}{\phi(x)}\,f(w)^{-1}$. & Accepted,
with one essential change: the twisted cocycle must be taken in the
centralizer $C_\phi$, not in $A$ (see item 1.8). The action is
\eqref{N-eq:gphi}; the formula for the restriction is
Lemma~\ref{N-lem:restrict}. The displayed formula in the report has a
misplaced parenthesis, but its conclusion is correct. \\
1.5 & In the definition of cohomologous $2$-cocycles, a prime is missing
on $\alpha'$. & Moot. The definition was replaced by Springer's
(Definition~\ref{N-def:H2}). \\
1.6 & The conjugation conventions $w^{-1}i(x)w=i(gx)$ and
$s(g)xs(g)^{-1}=gx$ are inconsistent; use left conjugation. & Accepted.
\eqref{N-eq:actN}. The report itself uses two conventions for $\sigma$
(once $s(g_1)s(g_2)s(g_1g_2)^{-1}$, once its inverse); Draft~1 uses the
first throughout. \\
1.7 & In the proof of semi-exactness, $w_i$ must be $s(g_i)$. & Moot. That
proof was replaced by Theorem~\ref{N-thm:main}, (b)$\Rightarrow$(c). \\
1.8 & The ``Conjecture'' (exactness at $H^1(E^\times,A)^\Gamma$) becomes a
theorem once the kernel of $\partial$ is defined as the set of classes with
neutral image in $H^2(\Gamma,A,\overline{\alpha})$; attributed to Springer,
Prop.~1.20. & \textbf{Not accepted as stated.} The two conditions on
$c(g)$ listed in the report are correct, and they are condition (c) of
Theorem~\ref{N-thm:main}; the formula $f'(i(x)s(g))=\phi(x)c(g)$ is
\eqref{N-eq:fc}. But these conditions are not equivalent to neutrality of
a class in $H^2(\Gamma,A,\overline{\alpha})$. Neutrality in $A$ allows the
elements $b_g$ to be changed by arbitrary elements of $A$, which destroys
the condition $b_g\in T_\phi(g)$. The correct target is
$H^2(\Gamma,C_\phi,\kappa_\phi)$. Example~\ref{N-ex:needed} gives
counterexamples to the statement with $A$ (the dihedral group of order
$8$ mapping to $Q_8$, and dihedral parameters in $\SL_2(\C)$). The
attribution to Springer's Prop.~1.20 could not be checked
(Section~\ref{sec:sources}) and is not used. \\
1.9 & In the conjecture, $f'(wi(x))=f'(w)f(gx)$ is not the cocycle
condition, which needs $\act{g}{f(x)}$, and a conjugation by $c(g)$ is
missing. & Accepted. The correct relations are \eqref{N-eq:fxw} and the
definition \eqref{N-eq:T} of $T_\phi(g)$. \\
2.1 & Add a remark on topology and continuity. & Accepted.
Proposition~\ref{N-prop:cont}. The relevant fact is that $E^\times$ is
open in $W_{E/F}$. \\
2.2 & Define neutral elements of $H^2$ and state exactness with respect
to them. & Accepted. Definition~\ref{N-def:H2},
Example~\ref{N-ex:H2}(c), Theorem~\ref{N-thm:main}(d). \\
3.1 & Convert from \AmS-\TeX\ to \LaTeX, remove unused macros and the old
affiliation. & Accepted. \\
3.2 & Misprints: ``p-adic'', ``pull-pack'', ``(semi-)exact'',
$\Aut(A)/A$ for $\Aut(A)/\Int(A)$, the broken sentence in the proof of
(1), ``an notions'', ``cocyles'', $r$ for $p$, ``$A$-groups'' for
``$\Gamma$-groups'', $g\in G$ for $g\in\Gamma$, ``cohomlogous'',
``particlar'', ``choosen''. & All accepted. The text was rewritten, so
most of these sentences no longer occur. \\
\bottomrule
\end{longtable}

\section{Comparison with the original manuscript}\label{sec:comparison}

\subsection{Statements of the original which were wrong or ill-posed}

These are recorded here and not in the note. Where a point is worth
keeping, the note presents it as a natural question followed by an
example.

\begin{enumerate}[itemsep=2pt]
\item \emph{Introduction.} ``$E^\times$ is isomorphic to $W_F/W_F^c$''
  should be $W_E/W_E^c$ (referee).
\item \emph{Introduction.} ``If $G$ is $E$-split but not $F$-split then
  $\Gamma$ acts on $\widehat{G}(\C)$ non-trivially'' is false for
  non-split inner forms of split groups. The action is trivial exactly
  when $G$ is an inner form of a split group. (Found in this revision.)
\item \emph{Lemma, hypothesis on $s$.} The condition
  $s(w^{-1}gw)=w^{-1}s(g)w$ is not meaningful and was removed (referee).
\item \emph{Lemma (1).} ``Non-trivial class'' needs $E\neq F$ (referee).
  The proof contained a broken sentence.
\item \emph{Claim under ``Restriction''.} The action $w^*f(x)=f(wx)$ is
  not defined on functions on $E^\times$, and the identification
  $H^1(E^\times,A)^\Gamma=\Hom_\Gamma(E^\times,A)/A$ is false (referee).
  In the note: the question after Example~\ref{N-ex:torus1} and
  Example~\ref{N-ex:toy-noeq}.
\item \emph{Definition of $\partial$.} The original put
  $\partial f=f\circ\sigma$ for $\Gamma$-equivariant $f$, with the given
  action $\alpha$, as a $2$-cocycle of $\Gamma$ in $A$. Its definition of
  a $2$-cocycle omitted Springer's condition
  $\alpha(g)\alpha(h)=\Int(u(g,h))\alpha(gh)$. With that condition,
  $(f\circ\sigma,\alpha)$ is a $2$-cocycle only when $f\circ\sigma$ is
  central in $A$. So the original $\partial$ is defined only for
  homomorphisms with $f\circ\sigma$ central. (Found in this revision.) The
  case of central $f$ is Corollary~\ref{N-cor:central}, where the original
  formula is correct up to the convention for $\sigma$.
\item \emph{Lemma (2), ``semi-exact at $H^1(E^\times,A)^\Gamma$''.} As a
  consequence of the previous two items the statement was not well posed.
  Its proof replaced $s(g_i)$ by arbitrary $w_i$ (referee). The correct
  statement is Theorem~\ref{N-thm:main}, (b)$\Rightarrow$(d).
\item \emph{Conjecture (exactness at $H^1(E^\times,A)^\Gamma$).} With
  target $H^2(\Gamma,A,\overline{\alpha})$ the conjecture is false, even
  when $\phi\circ\sigma$ is central: Example~\ref{N-ex:needed}. With
  target $H^2(\Gamma,C_\phi,\kappa_\phi)$ it is true:
  Theorem~\ref{N-thm:main}. For central $\phi$, and in particular for
  abelian $A$, the original conjecture is true as stated:
  Corollary~\ref{N-cor:central}. The proposed formula
  $f'(w)=c(g)f(ws(g)^{-1})$ has the factors in the wrong order; the
  correct one is \eqref{N-eq:fc}.
\end{enumerate}

\subsection{Original results and new results}

\begin{center}
\begin{tabular}{@{}p{0.5\textwidth}p{0.44\textwidth}@{}}
\toprule
Draft 1 & Status \\
\midrule
Lemma~\ref{N-lem:factor}(i)--(iii) & Original Lemma (1), corrected \\
Lemma~\ref{N-lem:factor}(iv), Example~\ref{N-ex:toy},
Example~\ref{N-ex:weil-quad} & New (standard) \\
Lemmas~\ref{N-lem:semidirect}, \ref{N-lem:ext-class}, \ref{N-lem:twist}
& New (standard) \\
Theorem~\ref{N-thm:exact} & Original Lemma (2), first three terms;
proofs kept, with corrections \\
Lemmas~\ref{N-lem:actions}, \ref{N-lem:restrict}, \ref{N-lem:T} & New \\
Definition~\ref{N-def:partial}, Proposition~\ref{N-prop:welldef},
Theorem~\ref{N-thm:main} & New; replaces the original connecting map and
settles the original conjecture \\
Corollary~\ref{N-cor:central} & Original formula for $\partial$, now
proved in the case where it applies \\
Proposition~\ref{N-prop:fiber} & New \\
Section~\ref{N-sec:graph} (Proposition~\ref{N-prop:derived},
Theorem~\ref{N-thm:graph}, Corollary~\ref{N-cor:cyclic}) & New \\
Section~\ref{N-sec:langlands}, Section~\ref{N-sec:questions}, glossary &
New \\
All examples & New \\
\bottomrule
\end{tabular}
\end{center}

\noindent ``New'' means new relative to the original manuscript. No claim
of novelty relative to the literature is made; see the paragraph ``Related
literature'' of the note.

\subsection{Changes of notation and convention}

\begin{itemize}[itemsep=1pt]
\item $\sigma(g,h)=s(g)s(h)s(gh)^{-1}$ replaces the original
  $s(gh)s(h)^{-1}s(g)^{-1}$ (its inverse). Noted in a footnote of the
  note.
\item Springer's conventions for $2$-cocycles are used. The original
  remarked that its $2$-cocycles differed from Springer's by an inversion.
\item The action of $g$ on $a$ is written $\act{g}{a}$ throughout; the
  original mixed $ga$ and $g(a)$.
\item Cohomologous: $f'(g)=a^{-1}f(g)\,\act{g}{a}$ (the original wrote the
  same relation with $f$ and $f'$ exchanged).
\item The general extension is written $1\to N\to W\to\Gamma\to1$; the
  original worked with $N=E^\times$ and remarked that $A$ could be any
  algebraic $\Gamma$-group over $\C$. Algebraicity is not used.
\item The original had a single unnumbered Lemma. Draft~1 numbers
  statements by section, in order of appearance.
\end{itemize}

\section{Verification}\label{sec:verification}

The script \texttt{verification/infres\_checks.py} treats twenty-seven
triples (extension $W$, normal subgroup $N$, $\Gamma$-group $A$), of which
twenty-two have $N$ abelian (twenty from Draft~1 and two added for
Draft~2: the counterexample $S_3$ of Corollary~\ref{N-cor:central} and the
Klein four-group example with $A=\mathbb{Z}/4$) and five have $N$
non-abelian ($N=Q_8$ in $Q_8\times\mathbb{Z}/2$ and in the central product
of $Q_8$ and $\mathbb{Z}/4$; $N=A_4$ in $S_4$). For non-abelian $N$ it
checks Lemma~\ref{N-lem:BD}, Theorem~\ref{N-thm:general}(a)--(c) and
Proposition~\ref{N-prop:fiber-general}. For abelian $N$ it also checks
that $B_\phi$ is a group of order $|C_\phi|\,|W|$ with $D_\phi$ normal.
The triples have
$|W|\le24$ and $|A|\le24$, with $\Gamma$ of order $2$, $3$, $4$ (Klein) and
$6$ (non-abelian), split and non-split, with trivial and non-trivial
actions. For each it enumerates all $1$-cocycles of $W$, $\Gamma$ and $N$
by brute force and checks: the factor set identity~\eqref{N-eq:2coc};
Theorem~\ref{N-thm:exact}; Lemma~\ref{N-lem:T}(ii);
Proposition~\ref{N-prop:welldef}(i),(ii); the equivalences (a), (b), (c),
(d) of Theorem~\ref{N-thm:main} and the bijection $f\leftrightarrow c$;
Proposition~\ref{N-prop:fiber}; and Theorem~\ref{N-thm:graph} for a
standard presentation of $\Gamma$. All checks pass. The output table also
confirms the counts in Examples~\ref{N-ex:toy-noeq}
and~\ref{N-ex:toy-obstruction}.

Not machine-checked: Proposition~\ref{N-prop:canonical} (the equation it
reduces to is checked), Propositions~\ref{N-prop:reconstruct}
and~\ref{N-prop:lifts}, Remark~\ref{N-rem:graph-nonab}, and the examples
over $p$-adic fields (Examples~\ref{N-ex:torus2},
\ref{N-ex:dihedral}, \ref{N-ex:unramified} and the running example). They
rest on facts (W1)--(W3) and on hand computation with $2\times2$ matrices.

\section{Sources consulted and not consulted}\label{sec:sources}

\begin{itemize}[itemsep=2pt]
\item \emph{Not consulted directly:} Springer \cite{Spr}. The definition
  of $2$-cocycles, equivalence and neutral classes in
  Definition~\ref{N-def:H2} is written from memory of the standard
  definition. The reference ``1.13--1.17'' is taken from the way the paper
  is cited in M.~Borovoi and D.~A.~Timashev, arXiv:2106.14871, \S5. The
  original manuscript's reference to Springer's Prop.~1.19 and the
  referee's reference to Prop.~1.20 were not checked and are not used. All
  statements about $H^2$ that the note needs are proved in the note
  (Lemmas~\ref{N-lem:ext-class} and~\ref{N-lem:twist}).
\item \emph{Not consulted directly:} Serre, \emph{Galois cohomology}. That
  \S I.5.8 contains the three-term sequence for a closed normal subgroup
  is taken from its citation in arXiv:2011.08773, \S4.1. The note cites
  ``Ch.~I, \S5''.
\item \emph{Not consulted directly, cited from memory:} Serre,
  \emph{Local fields} (Ch.~VII \S6 is the original manuscript's citation;
  Ch.~VIII \S4 for the cohomology of cyclic groups); Tate; Artin--Tate;
  Borel; Borovoi; Giraud; Hochschild--Serre; Langlands; Mac~Lane;
  Gross--Tucker; Zaslavsky; Baker--Norine. Section and chapter numbers for
  these should be checked before any submission.
\item \emph{Web search (one query, November 26, 2026)} on the non-abelian
  inflation-restriction sequence and neutral classes in Springer's $H^2$.
  No reference stating Theorem~\ref{N-thm:main} in this form was found,
  but the search was not exhaustive.
\end{itemize}



\section{Amendments to Draft 1}\label{sec:amend}

Draft~1 was amended on November 26, 2026, before it was merged into the
main branch, at the first author's request. The draft number was kept.
\begin{itemize}[itemsep=1pt]
\item \emph{Authorship.} Claude was added as second author. The first
  sentence of the abstract, the Acknowledgments and the sentence on
  related literature were reworded accordingly.
\item \emph{Glossary.} Rewritten and enlarged from 31 to 61 entries. Each
  entry is self-contained, and most now include an example, usually from
  the toy example $N=\Z/4$, $\Gamma=\Z/2$, $A=Q_8$. New entries:
  abelianization, central homomorphism, closed walk, coboundary,
  cohomology set $H^1$, covering map, defining relators, dihedral
  parameter, digraph, dual group, extension of groups, fiber, free group,
  holonomy, inner and outer automorphisms, inner form, monodromy, norm-one
  torus, pointed set, potential, quadratic character, reciprocity map,
  group of relations, split extension, switching, twisted norm, twisting,
  unramified, walk, Weil group.
\item No statement, proof or numbering in Sections~1--7 was changed by
  these two amendments.
\item \emph{Running example.} At the first author's request a detailed
  running example was added: the Weil group of the unramified quadratic
  extension $\mathbb{Q}_3(i)/\mathbb{Q}_3$, realized as
  $E^\times\sqcup E^\times s$ with $sxs^{-1}=\overline{x}$ and $s^2=3$,
  with coefficients $\C^\times$ (inversion action) and $\SL_2(\C)$
  (trivial action), and the tamely ramified characters $\chi_{k,z}$. It is
  in five parts, each placed at the end of a section so that no existing
  number changes: Examples~\ref{N-ex:run1}, \ref{N-ex:run2},
  \ref{N-ex:run3}, \ref{N-ex:run4} and~\ref{N-ex:run5}, with a new
  Figure~\ref{N-fig:run}. It computes the factor set, the effect of
  changing the section, the invariant classes, $C_\phi$, $T_\phi(\tau)$,
  $E_\phi$, $\partial\phi$, all cocycles in each fiber, the voltages, the
  derived graph, the holonomy and the potential. It also shows that the
  two toy extensions $D_4$ and $Q_8$ are quotients of this Weil group, and
  that the cocycle extending $\phi_{2,-1}$ is a parameter with image
  $Q_8$. The Introduction and its table were updated to mention it.
  All computations in the running example are by hand. The cases
  $k\in\{2,6\}$, $z=\pm1$ factor through the toy extensions, which are
  covered by the first two cases of the verification script; the rest is
  not machine-checked. The facts used about $\mathbb{Q}_3(i)^\times$ (its
  decomposition, the action of Frobenius on roots of unity, and
  $\mathrm{N}(1+3\mathcal{O})=1+3\mathbb{Z}_3$) are standard and were
  written from memory.
\end{itemize}

\section{Response to the second referee report (ChatGPT)}\label{sec:chatgpt}

The report was written from the first-pushed state of Draft~1, so its
remarks on the glossary and on the lack of non-cyclic and detailed
examples were partly answered already by the amendments of
Section~\ref{sec:amend}. Every mathematical claim of the report was
checked before use. All of them were found correct.

\begin{longtable}{@{}p{0.06\textwidth}p{0.36\textwidth}p{0.52\textwidth}@{}}
\toprule
Item & Referee's point & Response \\
\midrule
\endhead
M1 & ``The extension $W$ is encoded by the covering'' is too strong; give
the explicit reconstruction $W\cong(N\rtimes\mathcal{F})/\{(\nu(r)^{-1},r)\}$.
& Accepted. The sentence was replaced, with an example of two different
groups having the same covering, and
Proposition~\ref{N-prop:reconstruct} was added (with a converse, part
(ii), not in the report). \\
M2 & In Example~\ref{N-ex:torus1}, ``$\mathrm{res}$ is injective on the
fiber of the base point'' is awkward; for abelian coefficients
restriction is injective. & Accepted, reworded as proposed. \\
M3 & Corollary~\ref{N-cor:central}(ii) needs the hypothesis that $\phi$
is equivariant; counterexample $S_3$. & Accepted. This was an error of
Draft~1. The hypothesis was added, the counterexample is quoted in the
proof, and it was added to the verification script. \\
M4 & The harmonic morphism remark should specify the category of graphs;
Baker--Norine treat finite loopless multigraphs. & Accepted.
Remark~\ref{N-rem:harmonic} now notes that the Cayley digraphs used have
no loops (the generating set is assumed not to contain $1$) and states
the local definition for locally finite loopless multigraphs. \\
M5 & Distinguish the abstract-group $H^2$ from Springer's
Galois-cohomological one. & Accepted. A paragraph was added to
\S\ref{N-ss:scope}. \\
M6 & The statement on the triviality of the dual-group action should
refer to the pinned action. & Accepted, in \S1.1 and in a new glossary
entry. \\
\S5 & Choice-free obstruction extension
$\mathcal{O}_\phi=(W\times_\Gamma E_\phi)/D_\phi$, with
$[\mathcal{O}_\phi]=\partial\phi$. & Accepted and verified.
Proposition~\ref{N-prop:canonical}. \\
\S6 & The first open question of Draft~1 (non-abelian $N$) can be
answered: $\phi$ extends if and only if $B_\phi/D_\phi$ splits. &
Accepted and verified, including the report's argument with the inverse
image $L$ of a splitting. New Section~\ref{N-sec:nonabelian}:
Lemma~\ref{N-lem:BD}, Theorem~\ref{N-thm:general},
Remark~\ref{N-rem:compare}. Added beyond the report: the cocycle
equation~\eqref{N-eq:cgen} in this generality, the fibers
(Proposition~\ref{N-prop:fiber-general}),
Examples~\ref{N-ex:pauli}, \ref{N-ex:clifford}, \ref{N-ex:runH}, the
graph criterion for non-abelian $N$ (Remark~\ref{N-rem:graph-nonab}), and
machine checks. The question in Section~\ref{N-sec:questions} was
replaced. \\
7.1 & Emphasize that the split extension $D_4$ has a non-neutral
obstruction and the non-split $Q_8$ a neutral one. & Accepted; paragraph
added to Example~\ref{N-ex:toy-obstruction}. \\
7.2 & Display $\partial$ for the norm-one torus as a homomorphism to
$\{\pm1\}$. & Accepted, in Example~\ref{N-ex:torus2}. \\
7.3 & Add the determinant formula for an induced representation. &
Accepted, in Example~\ref{N-ex:dihedral}, with the direct check
$\det c=-\chi(a)$. The formula is quoted as well known, without a
reference. \\
7.4 & Compute the fiber in the unramified example for regular $t$:
$T[n]$, of order $n^r$. & Accepted, in Example~\ref{N-ex:unramified}. \\
8.1, 8.2 & Relation module, homology of the Cayley graph, cellular chain
sequence, presentation formula for $H^2$, dictionary table. & Accepted.
New \S\ref{N-ss:relmod}, with Proposition~\ref{N-prop:lifts} (proved),
the sequence~\eqref{N-eq:relseq} and formula~\eqref{N-eq:H2pres} (cited),
the table, and Example~\ref{N-ex:twocycle}. It replaces the remark ``The
universal voltage'' of Draft~1. The report states the chain sequence for
finite $\Gamma$ and finite generating sets; this restriction is not
needed and is not made. \\
8.3 & Add a non-cyclic example: $Q_8$ over the Klein four-group. &
Accepted. Example~\ref{N-ex:klein} and Figure~\ref{N-fig:klein}. \\
8.4 & Develop the question on harmonic morphisms that are not coverings
(graphs of groups, groupoid voltages). & Accepted as an expanded
Question~\ref{N-q:harmonic}; not developed further, as the report itself
suggests. \\
9.1 & Motivate the centralizer in the Introduction. & Accepted, in
\S1.1. \\
9.2 & Notation subsection. & Accepted, \S\ref{N-ss:notation}. \\
9.3 & Expand the glossary; proposed entries. & Mostly done already in the
amended Draft~1. Added now: augmentation ideal, derivation, fiber
product, graph homology, Klein four-group, obstruction extension, pinned
action, relation module. Not added: Baer sum and class formation, which
are not used in the note. \\
9.4 & Divide the literature review thematically; consider Brown and
Serre's \emph{Trees}. & Accepted. \\
9 & Diagram of the two extensions before the definition of the
obstruction class. & Accepted, \eqref{N-eq:diagram}. \\
9 & Modernize the font package. & Not done. The house style prescribes
Times through \texttt{mathptmx}. \\
10 & Reorganization with a new section on non-abelian kernels; optional
new title. & The organization was adopted. The title was kept; see
Section~\ref{sec:open}. \\
\bottomrule
\end{longtable}

\section{Changes from Draft 1 to Draft 2}\label{sec:d1d2}

\begin{itemize}[itemsep=1pt]
\item \emph{Error corrected.} Corollary~\ref{N-cor:central}(ii) lacked
  the hypothesis that $\phi$ is equivariant (item M3).
\item \emph{Statements qualified.} Items M1, M2, M4, M5, M6 above.
\item \emph{New results.} Proposition~\ref{N-prop:canonical};
  Section~\ref{N-sec:nonabelian} (Lemma~\ref{N-lem:BD},
  Theorem~\ref{N-thm:general},
  Proposition~\ref{N-prop:fiber-general});
  Propositions~\ref{N-prop:reconstruct} and~\ref{N-prop:lifts};
  Remark~\ref{N-rem:graph-nonab}. The first four were proposed by the
  referee; see the Acknowledgments of the note for the attribution.
\item \emph{New examples.} \ref{N-ex:pauli}, \ref{N-ex:clifford},
  \ref{N-ex:runH} (a new part of the running example, for a non-abelian
  normal subgroup of the Weil group of $\mathbb{Q}_3(i)/\mathbb{Q}_3$),
  \ref{N-ex:twocycle}, \ref{N-ex:klein}; additions to
  \ref{N-ex:torus2}, \ref{N-ex:toy-obstruction}, \ref{N-ex:dihedral},
  \ref{N-ex:unramified} and~\ref{N-ex:run3}. New Figure~\ref{N-fig:klein}.
\item \emph{Exposition.} Abstract and Introduction rewritten in part;
  notation table; diagram~\eqref{N-eq:diagram}; thematic literature
  paragraph; eight new glossary entries (69 in all); two new references
  (Brown; Serre, \emph{Trees}).
\item \emph{Authorship and acknowledgments.} ChatGPT is named in the
  abstract and its contributions are listed in the Acknowledgments, at the
  first author's request.
\item \emph{Assumption added.} The generating set $\overline{S}$ is
  assumed not to contain $1$.
\end{itemize}

\paragraph{Concordance of numbering.} Sections~1--3 keep their numbers,
except that \S1.4 (Notation) is new and the later subsections of
Section~1 shift by one. In Section~4, Proposition~\ref{N-prop:canonical}
is new, and Draft~1 numbers 4.6--4.13 become 4.7--4.14. Section~5 is new.
Draft~1 Sections 5, 6, 7 are now Sections 6, 7, 8. Within the graph
section: Draft~1 5.1--5.3 are now 6.1--6.3; 5.4 (harmonic morphisms) is
6.5; 5.5 (the universal voltage) was absorbed into \S\ref{N-ss:relmod};
5.6--5.10 are now 6.8--6.12; 5.11 (balance) is 6.15; 5.12 (running
example) is 6.16. Draft~1 6.1--6.4 are now 7.1--7.4, and Questions
7.1--7.4 are now 8.1--8.4, with a new Question~8.1. The running example
parts 4 and 5 of Draft~1 are now parts 5 and 6.

\section{Open items for the first author}\label{sec:open}

\begin{enumerate}[itemsep=1pt]
\item The year of the original manuscript is not known to the reviser.
  The date line reads ``Original draft: undated (1990s)''.
\item The bibliographic details listed in Section~\ref{sec:sources} should
  be checked against the sources. For Draft~2 this includes the
  references to Brown's \emph{Cohomology of groups} (Ch.~II, \S5 for the
  relation sequence; Ch.~IV for extensions) and to Serre's \emph{Trees},
  both cited from memory. The determinant formula for an induced
  representation in Example~\ref{N-ex:dihedral} has no reference.
\item Title. The second referee suggested, as an option,
  \emph{Inflation-restriction in non-abelian cohomology: centralizer
  obstructions, Weil groups, and voltage graphs}. The original title was
  kept.
\item Novelty. Neither referee nor the reviser made an exhaustive
  literature search. The note claims no novelty for
  Theorems~\ref{N-thm:main} and~\ref{N-thm:general}.
\item The house style sheet \texttt{joyner-math-paper-style-guide.md} is
  not in this repository. Draft~1 was written from the author's recorded
  style rulings: 12pt Times, letter paper, pdflatex, American spelling, AI
  disclosure at the head of the abstract, acknowledgments stating each
  AI's role, examples in every section, a glossary with self-contained
  entries, section-based numbering without lettered insertions, error
  discussion moved to this log.
\end{enumerate}

\begin{thebibliography}{9}
\bibitem[Spr]{Spr} T.~A.~Springer, Nonabelian $H^2$ in Galois cohomology,
Proc.\ Sympos.\ Pure Math.\ 9, Amer.\ Math.\ Soc., 1966, 164--182.
\end{thebibliography}

\end{document}
